% =============================================================================
% Section 2: The Anatomy of Innovation
% File: sections/02_anatomy.tex
% Compiles against shared main.tex preamble and macros.tex
% Reader question: What is innovation made of?
% Core points advanced: (2) Novelty is free; judgment is the inventor (E[N_tries]=1/q),
%                       (4) The skeptics are right about imitation and wrong about learning,
%                       (1) Innovation is a game, (3) Search = 1.
% =============================================================================

\section{The Anatomy of Innovation}\label{sec:anatomy}
% Reader question: What is innovation made of?
% Speaking points: 1 (Novelty is free), 2 (The four pillars), 3 (Taste & Eq. 1),
%                  4 (History vs. the referee), 5 (Where the pillars live)

What is innovation made of? Every human tradition of discovery rests on four pillars---aim, variation, judgment, and retention---where variation is free, judgment is costly, and taste (judgment learned in advance) sets how many tries a discovery takes.

\subsection{Novelty is free}\label{sub:anatomy-novelty}

Novelty was never the bottleneck in discovery; judgment was. In \emph{The Library of Babel}, \citet{borges1941babel} imagined a library containing every possible 410-page book of character permutations, and therefore the source code of every algorithm ever to be written. The library is useless to an inventor because nothing on its shelves tells the reader which volume works. A uniform random string generator is the most original machine there is, yet it invents nothing. In the standard research definition, creativity requires both originality and effectiveness \citep{runco2012standard}; in science and engineering, that effectiveness must be verified by an external test.

Controlled studies of automated research confirm that proposing unfamiliar ideas is far easier than judging which ones will work in practice. In a blinded evaluation with more than \num{100} natural language processing researchers, research ideas generated by language models were judged more novel than ideas written by human experts, scoring \num{5.83} versus \num{4.93} on a 10-point scale \citep{si2024ideas}. Yet when researchers carried out both sets of ideas, the generated proposals suffered a steep ideation--execution drop: their post-execution effectiveness scores fell by \num{1.87} points (from \num{6.00} to \num{4.13}), whereas human expert ideas held steady (from \num{4.83} to \num{4.78}), reversing the overall ranking in favor of human researchers \citep{si2025gap}. Unchecked variation produces plausible conjectures that lose ground once tested.

\begin{keyidea}
Novelty is free. Judgment is the inventor.
\end{keyidea}

\subsection{The four pillars of innovation}\label{sub:anatomy-pillars}

Every serious account of discovery names the same four functional pillars: \textbf{Aim} (Define), \textbf{Variation} (Propose), \textbf{Judgment} (Test), and \textbf{Retention} (Learn). \Cref{tab:traditions} shows these traditions agreeing from evolutionary biology to algorithmic self-play.

\begin{table}[t]
  \caption{\textbf{Every tradition names the same four pillars of innovation.} Across evolutionary biology, psychology, philosophy of science, engineering, and self-play, discovery couples a defined aim, a generator of candidate variations, an external referee, and a retention mechanism.}
  \label{tab:traditions}
  \centering
  \small
  \begin{tabular}{@{}l l l l l@{}}
    \toprule
    Tradition & Aim (Define) & Variation (Propose) & Judgment (Test) & Retention (Learn) \\
    \midrule
    \citet{darwin1859origin} & Fitness to an environment & Variation & Natural selection & Inheritance \\
    \citet{wallas1926art} & Preparation & Incubation, illumination & Verification & --- \\
    \citet{campbell1960bvsr} & A selection criterion & Blind variation & Selection & Selective retention \\
    \citet{popper1963conjectures} & The problem & Conjecture & Refutation & The surviving theory \\
    The scientific method & Question & Hypothesis & Experiment & Publication \\
    Engineering (DBTL) & Design target & Design and build & Test & Learn \\
    AlphaGo Zero \citep{silver2017mastering} & Win the game & Moves chosen with search & Game outcome & Network weights \\
    \textbf{\EZ{}} & \textbf{Problem encoding and metric} & \textbf{Program writer ($\writer$)} & \textbf{Referee ($\referee$) on hidden folds} & \textbf{Incumbent pool and weights} \\
    \bottomrule
  \end{tabular}
\end{table}

\textbf{Aim (Define)} specifies what problem to solve, how it is represented, and what counts as better. \citet{einstein1938evolution} noted that formulating a problem is often more essential than solving it, an observation echoed in P\'{o}lya's first rule to ``understand the problem'' \citep{polya1945solve}, Popper's problem situation \citep{popper1963conjectures}, and Hamming's question: ``What are the important problems of your field?'' \citep{hamming1986research}. In the innovation game, aim is encoded directly in the board.

\textbf{Variation (Propose)} produces candidate solutions. \citet[CP 5.171]{peirce1934cp} called this operation abduction, noting that it ``is the only logical operation which introduces any new idea.'' Candidates arise through recombinant growth over existing building blocks \citep{weitzman1998recombinant,arthur2009nature}, blind variation \citep{campbell1960bvsr}, bisociation across frames \citep{koestler1964act}, or combinational, exploratory, and transformational search within a conceptual space \citep{boden1990creative,boden1998creativity}. In \EZ{}, each variation is a complete executable program written by the program writer $\writer$.

\textbf{Judgment (Test)} evaluates a candidate against reality outside the proposer's mind. In science, judgment takes the form of attempted refutation \citep{popper1959logic,popper1963conjectures}, controlled experimentation, peer review, and the impartial referee of the Common Task Framework, which scores competing methods on a sealed test set \citep{donoho2017fifty}. In our game, the referee $\referee$ runs each proposed program and scores it on hidden folds.

\textbf{Retention (Learn)} keeps and builds on what survives. Human traditions store verified knowledge through biological inheritance \citep{darwin1859origin}, selective retention \citep{campbell1960bvsr,simonton2011bvsr}, and the cultural ratchet of textbooks and software libraries \citep{tomasello1999cultural,ellis2021dreamcoder}. An experiment, however, can do two jobs: evaluate an idea and improve the system that proposes the next idea. Archiving a winning program preserves the output; updating the generator's weights preserves the capability.

\begin{explainer}{Box 2: The Library of Babel and the Verification Gap}\label{box:babel}
Every algorithm that will ever be written already exists as a string in Borges's Library of Babel \citep{borges1941babel}. The challenge of invention is finding the rare string that generalizes. A random search reads the library book by book. An agent swarm equipped with an evaluator reads thousands of books per problem and keeps the best volume it finds, but starts over on the next problem. A player trained by self-play uses the referee's verdicts to learn which shelf to walk to---answering a new problem in one move (\searchone{}).
\end{explainer}

\subsection{Taste is judgment moved upstream}\label{sub:anatomy-taste}

When retention updates the generator itself, judgment moves upstream into variation. \citet[p.~324]{poincare1910creation} defined mathematical invention in exactly these terms: ``Invention is discernment, choice.'' We call this upstream discernment \emph{taste}.

Human expertise is taste stored in memory. Studying chess players from novices to grandmasters, \citet{degroot1965thought} and \citet{chase1973perception} showed that masters do not examine broader trees of moves than weaker players; they consider few moves, but the right ones. \citet[p.~155]{simon1992explanation} called this intuition ``nothing more and nothing less than recognition.'' Under time pressure, expert commanders and firefighters act on the first option they recognize \citep{klein1998sources,kahneman2011thinking}, embodying Polanyi's observation that ``we can know more than we can tell'' \citep[p.~4]{polanyi1966tacit}. AlphaGo Zero's policy network is taste made explicit in a neural network \citep{silver2017mastering}: through expert iteration \citep{anthony2017expert} and amortized optimization \citep{amos2023amortized}, it proposes moves that the referee will reward.

To formalize how taste sets the search budget, suppose that on a problem $p \in \Prob$, each try proposed by a generator beats the incumbent $\inc{p}$ with probability $q \in (0, 1]$. Under an idealized baseline of independent tries, a discovery requires $1/q$ tries on average:
\begin{equation}
  \mathbb{E}[N_{\text{tries}}] = \frac{1}{q}\,.
  \label{eq:taste}
\end{equation}
Here $N_{\text{tries}}$ is the number of candidate programs evaluated by the referee $\referee$ until the first win against $\inc{p}$, and $q$ is the per-try win probability.\footnote{\Cref{eq:taste} is an analytical idealization: in iterative agent loops or evolutionary search, tries are conditioned on prior outcomes and are not independent. We use $1/q$ to isolate the role of prior proposal quality at the first try ($B = 1$), and measure empirical multi-try curves directly in \cref{sec:results}.} \emph{Search length is the reciprocal of taste.}

\Cref{eq:taste} distinguishes four regimes on the plane of novelty and verified value (\cref{fig:novelty-value}). A random generator has $q$ near $0$: unbounded novelty without judgment. An imitating model has $q$ near $1$ for reproducing what humans already did, and $q$ near $0$ for beating the incumbent at the frontier. Agent swarms overcome low $q$ by paying for hundreds or thousands of tries $B$ at deployment. Self-play raises $q$ at the frontier by learning from the referee in advance, and \textbf{\searchone{} ($B = 1$) is the limit of high learned taste}. The same $q$ returns in \cref{sec:price} as the governing term of the economics: the expected cost per win is the cost per attempt divided by $q$ ($\Cwin = \Cattempt / q$).

\begin{figure}[t]
  \centering
  \includegraphics[width=\linewidth]{figures/fig2_novelty_value_schematic.pdf}
  \caption{\textbf{Novelty is free; judgment is the inventor (schematic).} Schematic plot of distance from what humans have done (horizontal) against value as judged by the referee $\referee$ (vertical). The solid curve marks the historical frontier. A random generator sits far right at near-zero value ($q \approx 0$); an imitating model reaches the historical frontier but does not cross it; agent swarms cross the frontier by spending hundreds to thousands of tries per problem; \EZ{} crosses the frontier in one try ($B = 1$) by learning taste from the referee during self-play.}
  \label{fig:novelty-value}
\end{figure}

\subsection{Why history bounds taste, and how a referee frees it}\label{sub:anatomy-history}

Why do human fields and imitating language models remain trapped near existing solutions? In human science, training instills the taste of the reigning paradigm \citep{kuhn1962structure}, and stepping stones that depart from consensus are routinely filtered out \citep{stanley2015greatness}. Analyzing \num{452} premature deaths of star life scientists, \citet{azoulay2019funeral} showed that subfields open to new investigators and highly cited outside ideas only after the incumbent leader dies---empirical evidence that science advances ``one funeral at a time.''

A language model trained to imitate human text inherits history's taste. That is the true content of the skeptics' argument. LeCun said in March 2025 that language models ``in their purest form, they are incapable of inventing new things'' \citep{lecun2025bigtech}; Wolf wrote in March 2025 that current training yields ``a country of yes-men on servers'' \citep{wolf2025einstein}; and Patel asked in August 2023 why models with ``basically the entire corpus of human knowledge memorized'' had not made ``a single new connection that has led to a discovery'' \citep{patel2023question}. Their critique is right about the training signal: next-token prediction on historical corpora rewards agreement with precedent. It is not a limit of neural networks.

The remedy is a teacher that is not history: a referee. AlphaGo Zero learned solely from game outcomes and discovered strategies beyond traditional Go canon \citep{silver2017mastering}. Earlier, AlphaGo estimated a one-in-ten-thousand chance that a human professional would play its Move 37 against Lee Sedol, yet the move helped win the game \citep{silver2016mastering,metz2016twomoves}. Moreover, discoveries made against a referee flow back to people: analyzing \num{5.8e6} professional moves across \num{58e9} counterfactual simulations, \citet{shin2023superhuman} showed that after AlphaGo, professional players' moves became both more novel and higher in quality. \emph{The skeptics are right about imitation and wrong about learning.}

\subsection{Where the pillars live}\label{sub:anatomy-where}

Where a system locates the four pillars determines whether it must rent judgment by the try on every problem or can own it in its weights. \Cref{tab:pillars-live} places human science, agent swarms, and \EZ{} side by side.

\begin{table}[t]
  \caption{\textbf{Swarms keep taste frozen in history; \EZ{} learns it from a referee.} Comparison of where the four pillars and taste reside across human science, agent swarms, and \EZ{}.}
  \label{tab:pillars-live}
  \centering
  \small
  \begin{tabular}{@{}l p{0.23\linewidth} p{0.31\linewidth} p{0.27\linewidth}@{}}
    \toprule
    & Human science & Agent swarms (\emph{autoresearch, AI Scientist, Eureka Machine}) & \EZ{} \\
    \midrule
    Aim & The researcher chooses & A human writes the goal & Problem encoding and metric, chosen by people for now \\
    Variation & The scientist's ideas & Agents propose and critique & The program writer $\writer$ makes one move \\
    Judgment & Experiments, peer review & Experiments; simulations and autonomous labs & The referee $\referee$ on hidden folds \\
    Retention & The literature & Logs, archives, a map of knowledge & Incumbent pool $\inc{p}$ and the writer's weights \\
    Where taste lives & Years of training in one head & Frozen in the pretrained prior: history's taste & Learned by self-play from the referee \\
    Tries per discovery & Many, slow & Hundreds per problem & One, on game day \\
    \bottomrule
  \end{tabular}
\end{table}

Richard Socher's \emph{The Eureka Machine} (October 2026) and Recursive organize automated science around four pillars that map directly onto ours \citep{socher2026eureka,recursive2026research}: ``a living map of all human knowledge'' serves retention; ``a unified model of physical reality,'' ``high-fidelity digital twins,'' and ``fully autonomous physical laboratories'' deliver faster, better judgment; and ``millions of AI agents'' that propose and critique ideas supply variation and filtering, paralleling agent loops such as \texttt{autoresearch} \citep{karpathy2026autoresearch} and the AI Scientist \citep{lu2024aiscientist,gottweis2025coscientist}.

Agent swarms and \EZ{} are complementary. Swarms scale test-time variation and build faster referees, while keeping taste frozen in the pretrained model's prior. \EZ{} asks what those referees should teach the policy itself. Every faster digital twin or autonomous laboratory built for agent swarms becomes a new board on which self-play can learn taste. In this work, people still supply the first pillar (\textbf{Aim}) by defining the problem encoding and metric; learning to choose problems is the next frontier (\cref{sec:meaning}). Each pillar already has a history in AI. The next section tells it, and finds the one piece no one has built.